\chapter{Immunology}

\medterm{Allergen}-- A substance capable of inducing an allergic (IgE-mediated) immune response in susceptible individuals. Allergens are typically proteins or glycoproteins with molecular weights between 10 and 70 kilodaltons, possessing structural features that promote IgE binding including protease resistance and the ability to crosslink IgE on mast cells. Major allergen families include profilins (pan-allergens in pollen, latex, fruits), lipid transfer proteins (LTPs, important in food allergy, particularly to peaches, apples, and nuts), and tropomyosins (shellfish allergens).

\medterm{Allergic Reactions}-- Hypersensitivity reactions mediated by IgE antibodies in response to otherwise harmless environmental substances (allergens). Type I immediate hypersensitivity reactions involve initial sensitization where dendritic cells present allergen to T helper 2 cells, stimulating B cell class switching to IgE production. IgE binds to high-affinity Fc-epsilon-RI receptors on mast cells and basophils, priming these cells. Upon re-exposure, allergen crosslinks surface-bound IgE, triggering degranulation and release of histamine, leukotrienes, and prostaglandins, causing vasodilation, bronchoconstriction, and increased vascular permeability.

\medterm{Allergic Rhinitis (Hay Fever)}-- Allergies occur when the immune system overreacts to harmless environmental substances such as pollen, dust mites, mold, or pet dander. Hay fever (allergic rhinitis) specifically involves nasal symptoms like sneezing, congestion, and itchy eyes triggered by airborne allergens. Management includes avoidance, antihistamines, nasal corticosteroids, and immunotherapy.

\medterm{Anaphylactic Shock}-- See Anaphylaxis.

\medterm{Ataxia-Telangiectasia}-- Autosomal recessive (ATM gene mutation, DNA damage response kinase). Progressive cerebellar ataxia (gait, truncal, speech), oculocutaneous telangiectasias (conjunctiva, ears, face, 3-6 years), sinopulmonary infections (IgA deficiency in 70%, IgG2 subclass deficiency), increased radiosensitivity, high malignancy risk (lymphoma, leukemia, 10-15%), endocrinopathies (diabetes, gonadal failure). Elevated AFP in 95%. Diagnosis: ATM sequencing, radiosensitivity assay. Treatment: IVIG for recurrent infections; avoid live vaccines and unnecessary radiation.

\medterm{Autoimmune Diseases}-- Autoimmune diseases occur when the immune system mistakenly attacks the body's own healthy tissues, producing chronic inflammation and organ damage. There are over 80 types, including rheumatoid arthritis, lupus, type 1 diabetes, multiple sclerosis, and Hashimoto's thyroiditis. Treatment typically focuses on suppressing immune activity with corticosteroids, immunomodulators, or biologics.

\medterm{Autoimmune Polyendocrine Syndromes}-- APS-1 (APECED, AIRE mutation): autosomal recessive, chronic mucocutaneous candidiasis, hypoparathyroidism, adrenal insufficiency (Addison's), plus type 1 diabetes, autoimmune thyroiditis, vitiligo, alopecia, pernicious anemia, hepatitis, asplenia. APS-2 (Schmidt syndrome): autoimmune adrenal insufficiency + autoimmune thyroid disease (Hashimoto's or Graves) + type 1 diabetes; associated with HLA-DR3/DR4. APS-3: autoimmune thyroid disease + type 1 diabetes (no adrenal). APS-4: other combinations of autoimmune endocrinopathies not fitting APS-1, 2, or 3.

\medterm{Autoimmunity: Mechanisms}-- Loss of tolerance: central (thymic negative selection defects, AIRE mutation -> APS-1) and peripheral (Treg dysfunction, anergy failure, ignorance loss). Molecular mimicry: cross-reactivity between pathogen and self (rheumatic fever: Strep M protein + cardiac myosin; Guillain-Barre: Campylobacter + GM1 ganglioside; reactive arthritis). Epitope spreading: immune response expands to new epitopes on same or different antigens (intramolecular/intermolecular). Bystander activation: inflammation activates autoreactive T cells non-specifically in the local milieu (e.g., in type 1 diabetes, myocarditis).

\medterm{Chediak-Higashi Syndrome}-- Autosomal recessive (LYST mutation, lysosomal trafficking regulator). Giant lysosomal granules in neutrophils, melanocytes, neurons. Partial oculocutaneous albinism, silvery hair, photophobia, nystagmus. Recurrent pyogenic infections (Staph, Strep), accelerated phase (HLH-like, 85%): fever, hepatosplenomegaly, pancytopenia, neurologic decline. Peripheral smear: giant azurophilic granules in neutrophils. Treatment: HSCT before accelerated phase; ascorbic acid improves chemotaxis.

\medterm{Chronic Granulomatous Disease}-- Chronic granulomatous disease (CGD) is a primary immunodeficiency caused by defects in the NADPH oxidase complex in phagocytes (neutrophils, monocytes, macrophages), impairing the respiratory burst and killing of catalase-positive organisms. The most common form is X-linked (CYBB gene mutation, encoding gp91phox, accounting for 65-70 percent of cases); autosomal recessive forms involve mutations in NCF1 (p47phox, 30 percent), NCF2 (p67phox), NCF4 (p40phox), or CYBA (p22phox). The defective phagocyte respiratory burst leads to recurrent infections with catalase-positive organisms including S. aureus, Serratia marcescens, Burkholderia cepacia, Nocardia, and Aspergillus.

\medterm{Chronic Mucocutaneous Candidiasis}-- Persistent/recurrent Candida infections of skin, nails, mucosa (oral, esophageal, vaginal) without systemic dissemination. Autosomal dominant (STAT1 gain-of-function) or recessive (IL17F, IL17RA, ACT1, RORC). Associated with autoimmunity (hypothyroidism, AIHA, ITP, alopecia), endocrinopathies, Staph infections. STAT1 GOF: multi-organ autoimmunity, cerebral aneurysms. Treatment: azole antifungals (fluconazole, itraconazole), JAK inhibitors (ruxolitinib) for STAT1 GOF; IL-17 pathway defects respond to recombinant IL-17 therapy.

\medterm{Common Variable Immunodeficiency}-- Common variable immunodeficiency (CVID) is the most common symptomatic primary immunodeficiency in adults, affecting approximately 1 in 25,000 individuals. It is characterized by hypogammaglobulinemia (low IgG with low IgA and/or IgM), impaired antibody response to vaccines, and increased susceptibility to infections. The pathogenesis is heterogeneous: approximately 50 percent of cases have identifiable genetic mutations (TNFSF13B/BAFF-R, TACI, MS4A1/CD20, IKZF1, CD19), while the remainder are idiopathic.

\medterm{Complement C1q Deficiency}-- Autosomal recessive (C1QA, C1QB, C1QC). C1q essential for classical pathway activation and immune complex clearance. 85-90% develop SLE-like disease (malar rash, photosensitivity, nephritis, anti-nuclear antibodies). Recurrent infections with encapsulated bacteria. Low CH50, low C1q, normal C3/C4. Treatment: SLE management (hydroxychloroquine, immunosuppressives); prophylactic antibiotics; IVIG for infections.

\medterm{Complement C2/C4 Deficiency}-- C2 deficiency: most common early complement deficiency (1:20,000 Caucasians). SLE-like disease (10-30%), recurrent sinopulmonary infections. C4 deficiency: similar, strong SLE association (C4A/C4B null alleles). Low CH50, low C2 or C4, normal C3. Treatment: SLE management, prophylactic antibiotics, IVIG for recurrent infections.

\medterm{Complement Deficiencies}-- Complement deficiencies are rare primary immunodeficiencies caused by genetic defects in complement pathway proteins, increasing susceptibility to specific infections and autoimmune diseases. Terminal complement deficiencies (C5-C9) impair membrane attack complex (MAC) formation, leading to recurrent Neisseria infections (meningococcal disease, gonococcemia) with a 10,000-fold increased risk compared to the general population. Vaccination against Neisseria meningitidis (quadrivalent ACWY and serogroup B vaccines), prophylactic antibiotics (penicillin), and prompt treatment of infections.

\medterm{Cytokine Release Syndrome}-- CRS: systemic inflammatory response from massive cytokine release (IL-6, IFN-gamma, TNF, IL-10) after T cell engaging therapies (CAR-T, bispecific antibodies, TGN1412). Grading (ASTCT): Grade 1 (fever), Grade 2 (hypoxia/hypotension responsive to O2/fluids), Grade 3 (hypoxia/hypotension requiring O2/vasopressors), Grade 4 (life-threatening). Management: tocilizumab (anti-IL-6R) first-line; corticosteroids for neurotoxicity (ICANS) or refractory CRS; anakinra for steroid-refractory; siltuximab (anti-IL-6) for refractory CRS.

\medterm{DiGeorge Syndrome}-- 22q11.2 deletion syndrome (most common microdeletion syndrome, 1:4000). Thymic hypoplasia/aplasia -> T-cell deficiency (variable, partial in 75%, complete in <1%), cardiac defects (conotruncal: TOF, truncus arteriosus, interrupted aortic arch), hypocalcemia (parathyroid hypoplasia), facial features, palatal abnormalities, learning disabilities. Immunodeficiency: variable CD4+ T-cell lymphopenia, impaired vaccine responses, autoimmunity risk. Treatment: calcium/vitamin D for hypocalcemia; IVIG if antibody deficiency is present; thymus transplant for complete DiGeorge.

\medterm{Factor H Deficiency}-- Autosomal recessive (CFH mutation). Factor H regulates alternative pathway (cofactor for factor I, accelerates decay). Uncontrolled alternative pathway activation -> C3 consumption -> low C3, normal C4. Atypical HUS (aHUS): microangiopathic hemolytic anemia, thrombocytopenia, acute kidney injury. C3 glomerulopathy (C3GN/DDD): C3-dominant deposits on IF. Treatment: eculizumab (anti-C5) for aHUS and C3G; plasma exchange for acute aHUS.

\medterm{Factor I Deficiency}-- Autosomal recessive (CFI mutation). Factor I cleaves C3b/C4b (cofactor: factor H, MCP, C4BP). Uncontrolled complement -> low C3, low C4, low factor H. Recurrent pyogenic infections, aHUS, C3 glomerulopathy, SLE-like disease. Treatment: eculizumab for aHUS/C3G; plasma infusion (contains factor I); IVIG.

\medterm{Granuloma}-- A distinctive pattern of chronic inflammation characterised by the aggregation of modified macrophages (epithelioid histiocytes), often with multinucleated giant cells (Langhans or foreign-body type), surrounded by lymphocytes. Granulomas form when macrophages cannot eliminate a persistent antigen or pathogen. Types: (1) Caseating (central necrosis): tuberculosis, fungal infections (histoplasmosis, coccidioidomycosis), syphilis. (2) Non-caseating: sarcoidosis, Crohn's disease, berylliosis, foreign body (talc, suture, silicone), granulomatosis with polyangiitis (GPA). (3) Suppurative (with neutrophils): actinomycosis, cat-scratch disease. Diagnosis: tissue biopsy (histology), special stains (AFB for TB, GMS for fungi), culture, PCR for specific pathogens. Treatment: depends on aetiology (anti-TB therapy, corticosteroids for sarcoidosis, surgical excision for foreign body).

\medterm{Griscelli Syndrome}-- Three types: Type 1 (MYO5A mutation): primary neurologic impairment, no immunodeficiency; Type 2 (RAB27A mutation): immunodeficiency + partial albinism, HLH risk; Type 3 (MLPH mutation): only partial albinism. Type 2: silvery hair, large melanosomes in melanocytes, defective cytotoxic granule release -> HLH. Treatment: HSCT for type 2; supportive for types 1 and 3.

\medterm{Hemophagocytic Lymphohistiocytosis}-- HLH: uncontrolled immune activation with cytokine storm (IFN-gamma, IL-6, IL-10, TNF). Familial HLH (FHL): PRF1 (FHL2), UNC13D (FHL3), STX11 (FHL4), STXBP2 (FHL5), RAB27A (FHL6/Griscelli type 2). Secondary HLH: infections (EBV, CMV), malignancy (lymphoma), autoimmune (SJIA/MAS, SLE), immunodeficiency. HLH-2004 criteria: 5/8 (fever, splenomegaly, cytopenias >=2 lineages, hypertriglyceridemia/hypofibrinogenemia, hemophagocytosis, low NK activity, ferritin >500, sCD25 >2400). Treatment: dexamethasone and etoposide (HLH-2004 protocol); emapalumab (anti-IFN-gamma) for refractory HLH; HSCT for familial HLH.

\medterm{Hermansky-Pudlak Syndrome}-- Autosomal recessive (10+ genes: HPS1-10, BLOC/AP-3 complexes). Oculocutaneous albinism, bleeding disorder (absent dense bodies in platelets -> prolonged bleeding time), pulmonary fibrosis (HPS-1, HPS-4), granulomatous colitis, neutropenia. Diagnosis: electron microscopy of platelets (absent dense bodies), genetic testing. Treatment: pulmonary fibrosis management (pirfenidone, nintedanib), avoid NSAIDs, HSCT not standard.

\medterm{HIV/AIDS Immunopathogenesis}-- HIV-1 infects CD4+ T cells, macrophages, dendritic cells via gp120 binding to CD4 + CCR5 (R5-tropic, early) or CXCR4 (X4-tropic, late). Reverse transcription -> integration -> transcription/translation -> budding. Acute retroviral syndrome (2-4 weeks): fever, lymphadenopathy, pharyngitis, rash, high viral load. Clinical latency (8-10 years): CD4 declines from 500-1500 to <200. AIDS: CD4 <200 or AIDS-defining OIs (PJP <200, toxoplasma encephalitis <100, CMV retinitis <50, MAC <50, cryptococcal meningitis <50).

\medterm{Hyper-IgM Syndromes}-- Hyper-IgM syndromes are a group of disorders characterized by normal or elevated IgM with low IgG, IgA, and IgE, due to defective class switch recombination (CSR) and somatic hypermutation (SHM). X-linked (CD40LG mutation, CD40 ligand deficiency, 70% of cases): T cells cannot provide CD40 signals to B cells; presents with recurrent pyogenic infections, Pneumocystis jirovecii pneumonia (PJP), cryptosporidiosis (sclerosing cholangitis), neutropenia, increased malignancy risk. Autosomal recessive: AID (activation-induced cytidine deaminase, AICDA), UNG (uracil N-glycosylase), or other CSR/SHM defects.

\medterm{Immune and Infectious Diseases}-- A medical field encompassing disorders of the immune system and illnesses caused by pathogenic microorganisms. Immune system disorders include autoimmune diseases (e.g., lupus, rheumatoid arthritis), immunodeficiencies, and hypersensitivity reactions. Infectious diseases range from viral respiratory infections to bacterial, fungal, and parasitic illnesses, requiring prevention through vaccination and treatment with antimicrobial agents.

\medterm{Immunoglobulin Replacement Therapy}-- Indications: primary antibody deficiencies (XLA, CVID, hyper-IgM, specific antibody deficiency), secondary hypogammaglobulinemia (CLL, myeloma, post-rituximab, post-transplant) with recurrent infections. IVIG: 400-600 mg/kg every 3-4 weeks (trough IgG >700-800 mg/dL); SCIG: 100-200 mg/kg/week (more stable levels, fewer systemic reactions). Products: varying IgA content (low IgA for anti-IgA antibodies), sucrose-free (renal safety), glycine, stabilizing agents. Adverse: headache, aseptic meningitis, infusion reactions, thrombosis; anaphylaxis in IgA deficiency (use low-IgA products).

\medterm{Immunotherapy and Biologic Treatments}-- Advanced therapies that harness or modify the immune system to treat diseases, particularly cancer, autoimmune disorders, and inflammatory conditions. Biologics are large-molecule drugs derived from living organisms that target specific immune pathways, such as TNF inhibitors, monoclonal antibodies, and checkpoint inhibitors. These treatments offer precision targeting but require careful monitoring for immune-related adverse effects.

\medterm{Inflammasomes}-- Multiprotein complexes activating caspase-1 -> cleavage of pro-IL-1beta/pro-IL-18 -> mature IL-1beta/IL-18 + pyroptosis (gasdermin D pores). Canonical: NLRP3 (crystals, ATP, K+ efflux, ROS, lysosomal rupture), NLRC4 (flagellin, T3SS), AIM2 (cytosolic DNA), Pyrin (RhoA inactivation). Non-canonical: caspase-11 (human caspase-4/5) activated by cytosolic LPS -> GSDMD cleavage -> pyroptosis + NLRP3 activation. Inflammasomopathies: CAPS (NLRP3 GOF), FMF (MEFV/pyrin LOF -> pyrin inflammasome overactivation), NLRC4 GOF (macrophage activation syndrome).

\medterm{Inflammation}-- The body's biological response to harmful stimuli such as pathogens, damaged cells, or irritants, characterized by redness, heat, swelling, pain, and loss of function. Acute inflammation is a protective process that eliminates the initial cause of injury and initiates tissue repair. Chronic inflammation, however, contributes to the pathogenesis of numerous diseases including atherosclerosis, diabetes, arthritis, and cancer.

\medterm{Innate Lymphoid Cells}-- ILCs: lymphoid lineage, no rearranged antigen receptors, mirror Th subsets. ILC1 (T-bet, IFN-gamma): intracellular pathogens, IBD. ILC2 (GATA3, IL-5/13): helminths, allergy, asthma, tissue repair. ILC3 (ROR-gammat, IL-17/22): extracellular bacteria/fungi, lymphoid organogenesis. LTi cells (ILC3 subset): lymph node/Peyer's patch development. NK cells (Eomes, perforin/granzyme): cytotoxic, missing-self recognition. ILCs in disease: ILC2 in asthma/CRSwNP (target: anti-IL-33, anti-TSLP, anti-IL-5); ILC3 in IBD and psoriasis; ILC1 in chronic inflammation.

\medterm{Interferon Signaling}-- Type I IFN (IFN-alpha/beta): all nucleated cells produce; receptor IFNAR (ubiquitous); JAK1/TYK2 -> STAT1/2/IRF9 (ISGF3) -> ISGs (hundreds: antiviral, antiproliferative, immunomodulatory). Type II IFN (IFN-gamma): T/NK cells; receptor IFNGR; JAK1/JAK2 -> STAT1 -> GAS elements. Type III IFN (IFN-lambda): mucosal epithelium; receptor IFNLR. IFN signatures: high in SLE (IFN-alpha), Aicardi-Goutieres (TREX1, RNASEH2, SAMHD1, ADAR1 mutations -> cytosolic nucleic acid accumulation -> IFN), SAVI (STING GOF), and COPA syndrome.

\medterm{Leukocyte Adhesion Deficiency}-- LAD-I (ITGB2 mutation, CD18 deficiency): absent CD11/CD18 integrins -> impaired neutrophil adhesion/transmigration -> delayed umbilical cord separation (>30 days), recurrent soft tissue infections without pus formation (leukocytosis >20,000), poor wound healing, periodontitis. Severe: death by age 2 without HSCT. LAD-II (SLC35C1 mutation, GDP-fucose transporter): absent sialyl Lewis X -> impaired rolling; mental retardation, short stature, Bombay blood group. LAD-III (FERMT3 mutation, kindlin-3): impaired integrin activation leading to a similar phenotype to LAD-I plus Glanzmann-like bleeding tendency.

\medterm{Lymph}-- A clear, watery fluid that circulates through the lymphatic system, derived from interstitial fluid that has entered the lymphatic capillaries. Lymph composition: water, electrolytes, small amounts of protein (especially in liver and intestine), white blood cells (predominantly lymphocytes, 95%, the remainder are monocytes, macrophages, and occasional granulocytes), cellular debris, microorganisms, and absorbed dietary lipids (chylomicrons, giving lymph a milky appearance--chyle). Lymph is formed when hydrostatic and osmotic pressures across the capillary wall cause net filtration of fluid into the interstitial space. Approximately 2-4 L of lymph per day drains through the thoracic duct into the left subclavian vein. Lymph flow is propelled by: contraction of smooth muscle in the wall of collecting lymphatics (intrinsic contractility), compression by skeletal muscle contraction, respiratory motion, and pulsation from adjacent arteries. One-way valves prevent backflow. Lymph nodes filter lymph, removing pathogens and antigens, and serving as sites for immune surveillance and activation. The lymphatic system has crucial roles in: fluid homeostasis (returns excess interstitial fluid to the bloodstream), immune surveillance (lymphocytes recirculate through lymph nodes, detecting and responding to antigens), and fat absorption (intestinal lacteals absorb dietary lipids).

\medterm{Lymph Node / Gland}-- Small, bean-shaped organs of the lymphatic system that filter lymph and serve as sites for immune cell activation. Lymph nodes are distributed throughout the body (approximately 500-600 in total) and are concentrated in regions: cervical, axillary, inguinal, mediastinal, mesenteric, para-aortic, and iliac chains. Structure: fibrous capsule, cortex (B-cell follicles--primary and secondary with germinal centres), paracortex (T-cell zone, high endothelial venules), and medulla (medullary cords with plasma cells and macrophages, medullary sinuses). The subcapsular sinus receives afferent lymphatic vessels. The hilum contains the efferent lymphatic vessel and blood vessels. Functions: filtration--macrophages in the sinuses phagocytose particulate matter and microorganisms; immune surveillance--antigens are presented to lymphocytes by dendritic cells, leading to B-cell and T-cell activation, antibody production, and generation of memory cells. Lymph nodes enlarge (lymphadenopathy) in response to infection, inflammation, or malignancy. Palpable lymph nodes in adults: >1 cm in most regions is considered abnormal, except inguinal nodes (1.5 cm). Sentinel lymph node: the first lymph node draining a tumour, biopsied to stage cancer spread (e.g., melanoma, breast cancer).

\medterm{Lymphadenopathy}-- Enlargement of lymph nodes (greater than 1 cm in diameter, or >1.5 cm for inguinal nodes). Lymphadenopathy can be localised (single or contiguous region) or generalised (three or more non-contiguous regions). Localised causes: infection (most common--pharyngitis/tonsillitis causing cervical adenopathy, wound infection causing axillary/inguinal adenopathy, TB causing supraclavicular/scalene adenopathy, cat-scratch disease causing axillary/epitrochlear adenopathy), malignancy (lymphoma, metastatic cancer, leukaemia), autoimmune (sarcoidosis, SLE, rheumatoid arthritis). Generalised causes: infectious (EBV, CMV, HIV, TB, toxoplasmosis, brucellosis), malignant (lymphoma, CLL, ALL), autoimmune (SLE, sarcoidosis, Still disease), drug-induced (phenytoin, allopurinol), and storage diseases (Gaucher, Niemann-Pick). Diagnostic evaluation: history (duration, size, tenderness, fever, night sweats, weight loss, exposure, travel), physical exam (location, size, consistency--hard/fixed suggests malignancy, soft/mobile suggests infection, matted suggests TB/sarcoidosis, fluctuant suggests abscess), and workup: full blood count, CRP/ESR, serology (EBV, CMV, HIV, toxoplasma), chest X-ray, ultrasound of nodes, and excisional biopsy (indications: node >2 cm, supraclavicular location, firm/fixed, no regression after 4-6 weeks, systemic symptoms, or abnormal CXR).

\medterm{Lymphangioma}-- A benign congenital malformation of the lymphatic system, composed of dilated lymphatic channels. Lymphangiomas are typically present at birth (90% diagnosed by age 2) and result from failure of the lymphatic system to develop normal connections with the venous system. Types: (1) Cystic hygroma: most common (50-60%), large multiloculated cysts, typically in the neck (posterior triangle) or axilla. Antenatal ultrasound may detect as a nuchal translucency (associated with Turner syndrome, Noonan syndrome, trisomy 21). (2) Cavernous lymphangioma: deep, involves neck, tongue, lip (macrodontia, macroglossia). (3) Lymphangioma circumscriptum: cutaneous vesicles on skin surface. Symptoms: soft, fluctuant, transilluminating mass, often painless. May enlarge rapidly with upper respiratory infection or trauma due to increased lymph production. Complications: infection (lymphangitis), intralesional bleeding, compression of airways, dysphagia, cosmetic deformity. Diagnosis: clinical, ultrasound (cystic, anechoic, septated), MRI (T2 hyperintense, fluid-fluid levels, no enhancement). Management: observation (30% may regress spontaneously, especially small ones), surgical excision (complete excision is curative, but recurrence risk of 10-40% if incomplete), sclerotherapy (OK-432, bleomycin, doxycycline, alcoholic solution of zein--Ethibloc, especially for cystic hygroma), sirolimus (mTOR inhibitor, for complex/refractory cases).

\medterm{Lymphedema (Lymphoedema)}-- Chronic swelling of a body part due to impaired lymphatic drainage, leading to accumulation of protein-rich interstitial fluid. Primary lymphedema: caused by congenital or hereditary lymphatic maldevelopment. Milroy disease (autosomal dominant, VEGFR3 mutation, congenital, present at birth, 2:1 female:male, lower limbs, non-pitting), Meige disease (lymphedema praecox, onset at puberty, 2:1 female:male, lower limbs up to knee, mild), and late-onset lymphedema (tarda, >35 years). Secondary lymphedema: more common, acquired damage to lymphatic vessels or nodes. Causes: cancer-related (axillary lymph node dissection for breast cancer--most common cause of upper limb lymphedema, pelvic node dissection for gynaecological/ urological cancers, melanoma), radiation therapy (lymphatic fibrosis), filariasis (Wuchereria bancrofti, Brugia malayi--most common cause worldwide, mosquito-borne nematode infection, tropical regions, causes massive leg and genital swelling, elephantiasis), trauma/surgery (lymph node excision), recurrent cellulitis, liposuction (may worsen), and venous disease (phlebolymphedema). Pathophysiology: lymphatic obstruction leads to accumulation of protein-rich fluid, which triggers chronic inflammation, adipose deposition, and fibrosis. Over time, the skin becomes thickened, hyperkeratotic, and papillomatous. Diagnosis: clinical (non-pitting oedema Stemmer sign--inability to pinch the skin of the dorsum of the second toe, specific for lymphedema), lymphoscintigraphy (radionuclide imaging of lymphatic transport, gold standard), MRI (honeycomb pattern, dermal thickening, oedema fluid), ICG fluorescence lymphangiography. Management: complex decongestive therapy (CDT, gold standard)--phase 1 (intensive, 2-4 weeks, manual lymphatic drainage MLD, compression bandaging, exercise, skincare) and phase 2 (maintenance, daytime compression garments, nighttime bandaging, self-MLD). Additional: pneumatic compression devices, laser therapy, liposuction (for non-pitting, fibrotic, late-stage lymphedema), surgical (lymphaticovenous anastomosis LVA, vascularised lymph node transfer VLNT, Charles procedure for severe elephantiasis). Antibiotic prophylaxis (penicillin V, 250 mg BD) for recurrent cellulitis. Long-term management with patient education, weight control, and skincare.

\medterm{Lymphocyte}-- A type of white blood cell (leukocyte) that is the primary cellular component of the adaptive immune system. Lymphocytes account for 20-40% of circulating white blood cells (1.0--4.8 $\times$ 10$^9$/L). Three main types: (1) T lymphocytes (T cells): develop in the thymus, responsible for cell-mediated immunity. Subsets: CD4+ helper T cells (Th1, Th2, Th17, Tfh, Treg) orchestrate immune responses by secreting cytokines; CD8+ cytotoxic T cells kill virus-infected cells, tumour cells, and allograft cells. (2) B lymphocytes (B cells): develop in the bone marrow, responsible for humoral immunity. Upon activation, B cells differentiate into plasma cells (secrete antibodies--immunoglobulins IgG, IgA, IgM, IgE, IgD) and memory B cells (long-lived, rapid response upon re-exposure). (3) Natural killer (NK) cells: part of the innate immune system, kill virus-infected cells and tumour cells without prior sensitisation (no antigen receptor rearrangement). Lymphocyte development: haematopoietic stem cells in bone marrow differentiate into common lymphoid progenitors. B cells mature in the bone marrow, T cells migrate to the thymus for maturation (positive selection for self-MHC recognition, negative selection against self-reactivity). After encountering antigen, lymphocytes undergo clonal expansion and differentiate into effector and memory cells. Lymphocyte recirculation: naive lymphocytes travel from blood to lymph nodes via high endothelial venules (HEVs), then to lymph, and back to blood. Lymphocyte disorders: lymphocytosis (viral infection, CLL), lymphopenia (HIV, chemotherapy, immunodeficiency), lymphoma (malignant proliferation of lymphocytes), and autoimmune diseases (dysregulated lymphocyte reactivity).

\medterm{Macrophage Activation Syndrome}-- MAS is a form of secondary HLH complicating rheumatic diseases (SJIA, Kawasaki, SLE, AOSD). MAS in SJIA: cytopenias, hyperferritinemia, elevated LDH, elevated triglycerides, low fibrinogen, transaminitis. High mortality (20-50%). IL-1/IL-6 driven. Treatment: high-dose corticosteroids, anakinra (IL-1 blockade), tocilizumab (IL-6 blockade), cyclosporine A, etoposide for severe.

\medterm{Mannose-Binding Lectin Deficiency}-- MBL2 polymorphisms (common, 10-30% heterozygotes, 1-5% homozygotes). MBL activates lectin pathway. Low MBL associated with increased infections in immunocompromised (chemotherapy, transplant), recurrent infections in children, autoimmune disease susceptibility. Diagnosis: serum MBL level, MBL2 genotyping. Treatment: usually asymptomatic; IVIG for recurrent infections in immunocompromised.

\medterm{Neutrophil Extracellular Traps}-- NETs: web-like structures of chromatin + granular proteins (MPO, NE, histone) released by neutrophils (NETosis: suicidal vs vital). Functions: trap/kill pathogens (bacteria, fungi, viruses); pathological in thrombosis (NETs provide scaffold for platelets, activate coagulation), autoimmunity (SLE: NETs expose LL-37/DNA -> pDC activation -> IFN-alpha; RA: citrullinated antigens in NETs -> ACPA), cancer (NETs trap circulating tumor cells -> metastasis), COVID-19 (immunothrombosis). DNase I degrades NETs; therapeutic DNase (dornase alfa) used in NET-related pathology.

\medterm{Pattern Recognition Receptors}-- PRRs recognize PAMPs (pathogen) and DAMPs (damage). TLRs: TLR1/2 (lipoproteins), TLR2/6 (lipoteichoic acid), TLR3 (dsRNA), TLR4 (LPS, MD-2, CD14), TLR5 (flagellin), TLR7/8 (ssRNA), TLR9 (CpG DNA). Cytosolic: RIG-I/MDA5 (RNA), cGAS (DNA -> cGAMP -> STING -> IFN), NOD1/2 (peptidoglycan). CLRs: Dectin-1 (beta-glucan), Mincle (mycobacterial TDM), DC-SIGN (mannose). NLRs: NOD1/2 (cytosolic bacteria), NLRP3 (inflammasome). Signaling: MyD88 (all TLRs except TLR3) -> NF-kB/IRF; TRIF (TLR3/4) -> IRF3/NF-kB activation.

\medterm{Phagocytosis}-- The process by which specialised cells (phagocytes) engulf and destroy foreign particles, microorganisms, and cellular debris. Professional phagocytes: neutrophils (most abundant, first responders, short-lived), macrophages (tissue-resident—Kupffer cells in liver, alveolar macrophages in lung, microglia in brain, osteoclasts in bone), dendritic cells (antigen presentation to T cells), and eosinophils (parasite defence). Phagocytosis steps: (1) Chemotaxis—phagocytes migrate to the site of infection guided by chemoattractants (IL-8, C5a, bacterial peptides). (2) Recognition and opsonisation—pathogens are coated with opsonins (IgG antibodies, complement C3b, collectins) that bind phagocyte receptors (Fc\(\gamma\) receptors, complement receptors, mannose receptors). (3) Engulfment—the phagocyte extends pseudopods around the particle, forming a phagosome. (4) Phagosome–lysosome fusion—the phagosome fuses with lysosomes to form a phagolysosome; the pathogen is killed by reactive oxygen species (respiratory burst—NADPH oxidase produces superoxide), nitric oxide, antimicrobial peptides (defensins, cathelicidins), and hydrolytic enzymes (lysozyme, proteases). (5) Exocytosis—degraded material is expelled. Defects in phagocytosis cause recurrent infections: chronic granulomatous disease (NADPH oxidase deficiency—inability to produce reactive oxygen species), Chediak–Higashi syndrome (defective phagolysosome fusion), and leukocyte adhesion deficiency (impaired migration).

\medterm{Properdin Deficiency}-- X-linked recessive (CFP gene). Properdin stabilizes alternative pathway C3 convertase. Recurrent Neisseria infections (meningococcemia, gonococcemia). Low AH50, normal CH50, low properdin. Treatment: meningococcal vaccination (ACWY + B), prophylactic penicillin, IVIG for severe infections.

\medterm{Regulatory T Cells}-- Tregs (CD4+ CD25+ FOXP3+): essential for peripheral tolerance. Thymic (tTreg) vs peripheral (pTreg) origin. FOXP3: master transcription factor; mutations cause IPEX. Mechanisms: IL-2 consumption (CD25 high), CTLA-4 (trans-endocytosis of CD80/86), TGF-beta/IL-10 secretion, cAMP transfer, granzyme/perforin killing of APCs, metabolic disruption (CD39/CD73 -> adenosine). Treg defects: IPEX (FOXP3), autoimmune polyendocrinopathy. Therapeutic: low-dose IL-2 expands Tregs (trials in SLE, GVHD); Treg adoptive transfer (trials in GVHD, type 1 diabetes, transplant rejection).

\medterm{Secondary Immunodeficiencies}-- Acquired immune defects from: HIV/AIDS (CD4+ T-cell depletion), malignancies (CLL: hypogammaglobulinemia; myeloma: hypogammaglobulinemia + impaired specific antibody; lymphoma), immunosuppressive medications (corticosteroids, chemotherapy, biologics: rituximab -> prolonged B-cell depletion and hypogammaglobulinemia; TNF inhibitors -> TB reactivation; alemtuzumab -> profound T-cell depletion), protein-losing states (nephrotic syndrome, enteropathy, burns), splenectomy (encapsulated bacteria risk, asplenia), and advanced age (immunosenescence).

\medterm{Severe Combined Immunodeficiency}-- Severe combined immunodeficiency (SCID) represents a group of genetic disorders characterized by profoundly defective T cell development and function, with variable B and NK cell involvement, resulting in severe susceptibility to life-threatening infections in infancy. It is the most serious primary immunodeficiency, with a combined incidence of approximately 1 in 50,000-100,000 live births. The classic form (X-linked SCID, accounting for approximately 45-50 percent of cases) is caused by mutations in the IL2RG gene (common gamma chain), resulting in absent T cells and NK cells with normal or elevated B cells (T-B+NK- SCID).

\medterm{T-Cells (T-Lymphocytes)}-- A type of lymphocyte that matures in the thymus and is the central effector of cell-mediated immunity. T-cells express the T-cell receptor (TCR) on their surface, which recognises peptide antigens presented by major histocompatibility complex (MHC) molecules on antigen-presenting cells. T-cells are divided into subsets by function and surface markers: (1) CD4+ helper T cells (Th)—recognise MHC class II, orchestrate immune responses. Subsets: Th1 (IFN-\(\gamma\), cell-mediated immunity against intracellular pathogens), Th2 (IL-4, IL-5, IL-13, humoral immunity, allergic responses), Th17 (IL-17, neutrophil recruitment, antifungal/extracellular bacterial immunity), Tfh (follicular helper T-cells—help B cells in germinal centres, antibody production), and Treg (regulatory T-cells, FOXP3+, suppress immune responses, maintain self-tolerance). (2) CD8+ cytotoxic T cells (CTL)—recognise MHC class I, kill virus-infected cells, tumour cells, and allograft cells via perforin/granzyme and Fas/FasL pathways. (3) \(\gamma\delta\) T-cells—innate-like, recognise non-peptide antigens, bridge innate and adaptive immunity. T-cell development: bone marrow progenitor \(\rightarrow\) thymus (positive selection in the cortex, negative selection in the medulla) \(\rightarrow\) naive T-cell \(\rightarrow\) circulation. Activation: requires TCR binding (signal 1) and co-stimulation (CD28–CD80/86, signal 2). Deficient T-cell function causes severe opportunistic infections (HIV/AIDS, severe combined immunodeficiency—SCID). CD4 count is the key marker of immune function in HIV (normal 500--1500 cells/\(\mu\)L; <200 defines AIDS).

\medterm{Th1/Th2/Th17/Tfh/Treg Balance}-- CD4+ T helper subsets defined by master transcription factors and signature cytokines: Th1 (T-bet, IFN-gamma, IL-12, STAT4): intracellular pathogens, autoimmunity (RA, MS, type 1 diabetes). Th2 (GATA3, IL-4/5/13, IL-4/STAT6): helminths, allergy, asthma. Th17 (ROR-gammat, IL-17A/F, IL-21, IL-22, IL-6+TGF-beta/IL-23/STAT3): extracellular bacteria/fungi, autoimmunity (PSA, AS, MS, IBD). Tfh (Bcl-6, IL-21, CXCR5): germinal center help, antibody affinity maturation, memory B cells. Treg (FOXP3, TGF-beta, IL-10, IL-35): suppression, immune regulation, tolerance.

\medterm{Trained Immunity}-- Innate immune memory: epigenetic/metabolic reprogramming of monocytes/macrophages/NK cells after initial stimulus (BCG, beta-glucan, Candida, oxidized LDL) -> enhanced response to secondary challenge (heterologous protection). Mechanisms: histone modifications (H3K4me3, H3K27ac), metabolic shifts (aerobic glycolysis, glutaminolysis, cholesterol synthesis), mTOR-HIF1alpha pathway. BCG vaccination protects against unrelated infections (neonatal sepsis, respiratory infections). Relevance: sepsis-induced immunoparalysis (trained immunity reversal).

\medterm{Transplant Immunology}-- Hyperacute rejection: preformed anti-donor antibodies (ABO, HLA) -> immediate graft failure (minutes-hours); prevented by crossmatch. Acute cellular rejection: recipient T cells recognize donor alloantigens (direct: donor APC -> recipient T cell; indirect: recipient APC presents donor peptides); days-months; treated with steroids, T cell depletion (thymoglobulin, alemtuzumab). Acute antibody-mediated rejection: de novo DSA -> C4d deposition, capillaritis, glomerulitis; treated with plasmapheresis, IVIG, rituximab, and bortezomib. Chronic rejection: interstitial fibrosis and tubular atrophy (IFTA) with DSA; poorly responsive to therapy.

\medterm{Vaccine Immunology}-- Live attenuated (MMR, varicella, rotavirus, yellow fever, oral typhoid, BCG, intranasal flu): replicate in host -> strong humoral + cellular immunity; contraindicated in immunocompromised. Inactivated/killed (polio IPV, hepatitis A, rabies, injectable flu): safer, weaker immunity, need boosters. Subunit/recombinant (hepatitis B, HPV, pertussis acellular, MenB, RSV): specific antigens, safe, need adjuvants. Toxoid (diphtheria, tetanus): inactivated toxin. mRNA (COVID-19 Pfizer/Moderna): lipid nanoparticles, highly effective, rapid development, require cold chain storage.

\medterm{Vaccines and Immunizations}-- Vaccines and immunizations are biological preparations that stimulate the immune system to develop protection against specific infectious diseases. By exposing the body to harmless components of a pathogen, vaccines enable the immune system to generate memory B and T cells that can mount a rapid response upon future exposure. Widespread immunization has dramatically reduced the burden of diseases such as polio, measles, influenza, and COVID-19, and remains one of the most cost-effective public health interventions for preventing morbidity and mortality.

\medterm{Wiskott-Aldrich Syndrome}-- X-linked recessive disorder (WAS gene mutation, encodes WASP protein, regulates actin polymerization in hematopoietic cells). Triad: microthrombocytopenia (small platelets, <70 fl), eczema, recurrent infections (encapsulated bacteria, viruses). Combined immunodeficiency with progressive T-cell lymphopenia. Increased risk of autoimmune disease (AIHA, ITP, vasculitis, IBD) and malignancy (lymphoma, 10-15%). Diagnosis: low WASP expression, WAS gene sequencing. Treatment: IVIG for infections, platelet transfusions for bleeding; splenectomy for refractory thrombocytopenia; HSCT as definitive cure.

\medterm{X-Linked Agammaglobulinemia}-- X-linked agammaglobulinemia (XLA, Bruton agammaglobulinemia) is caused by mutations in the BTK gene (Bruton tyrosine kinase), essential for B-cell development and maturation. Results in absent or very low mature B cells (<1% CD19+ cells), profoundly low all immunoglobulin isotypes, and absent antibody responses. Presents in male infants at 4-8 months (after maternal IgG wanes) with recurrent sinopulmonary infections (S. pneumoniae, H. influenzae, M. catarrhalis), enteroviral infections (chronic enteroviral meningoencephalitis), and lack of lymph nodes and tonsils. Diagnosis: low immunoglobulins, absent B cells, BTK sequencing. Treatment: lifelong IVIG or SCIG replacement therapy; avoid live vaccines; gene therapy is experimental.
